\chapter*{Preface}
\addcontentsline{toc}{chapter}{Preface}

Computational environmental science increasingly depends on the orchestration
of heterogeneous resources: numerical models compiled for HPC clusters, data
preprocessing and post-processing scripts, container images that freeze a
fragile scientific software stack, cloud instances rented for a single
forecast cycle, and --- more recently --- AI services consumed through web
APIs for summarization, anomaly detection, or decision support. The challenge
this book addresses is not how to write any one of these components, but how
to make their \emph{composition} explicit, reproducible, portable, and
auditable.

\emph{DAGonStar} answers this challenge with a deliberately small abstraction:
a workflow is a directed acyclic graph (DAG) of tasks, and a task is a Python
object that executes a command --- local, remote, containerized,
scheduler-submitted, an HTTP request, an LLM completion, or an importable
Python function --- in a particular environment. This minimalism is the
engine's central design virtue. It keeps workflow authoring close to ordinary
Python, while still providing the scientific guarantees that ad hoc shell
scripts cannot: explicit dependency graphs, automatic data staging, cycle
detection, checkpointing, and a uniform mechanism --- the \texttt{workflow://}
schema --- for expressing that one task consumes a file produced by another.

\paragraph{How this book is organized.}
The chapters follow, almost one-to-one, the project's own fourteen-lesson
tutorial sequence (\texttt{docs/tutorial/}), because that sequence was already
designed as an incremental, self-verifying curriculum: each lesson states an
objective, a concept, a complete code listing, the expected behavior, a
verification procedure, and a ``scientific practice note'' connecting the
mechanism to good research practice. This book preserves that pedagogical
backbone but expands every lesson into a fuller theoretical discussion,
situates it within the engine's architecture, and illustrates it with
original figures.

Part~I (Chapters~1--8) builds the conceptual and mechanical foundations:
scientific workflow theory, the DAGonStar architecture, installation, the
\texttt{workflow://} schema, configuration, and Lessons~1--7 (local workflows
through data staging). Part~II (Chapters~9--14) covers heterogeneous execution
backends and the engine's AI--HPC integration layer: Lessons~8--14 (Docker,
remote/Slurm execution, meta-workflows, asynchronous launch, LLM tasks, Web
tasks, and Native Python tasks), plus a dedicated chapter on wrapping external
scientific software. Part~III (Chapters~15--17) turns to computational
environmental science proper, using the repository's environmental
application example as an extended case study and closing with a synthesis
chapter that connects every preceding mechanism into the shape of an
operational, AI-augmented forecasting chain. Appendices~A--D collect reference
material --- the public API, developer practices, the full configuration
model, and checkpoint semantics --- intended for use after the narrative
chapters, while working with real workflows.

\paragraph{Intended audience.}
This book is written for graduate students, researchers, and research software
engineers in computational and environmental science who need to compose
heterogeneous computational pipelines reproducibly. Working knowledge of
Python and of basic Unix/HPC concepts (shells, SSH, job schedulers) is
assumed; no prior exposure to workflow engines is required.

\paragraph{A note on fidelity to source material.}
Every code listing in this book is reproduced, without semantic alteration,
from the DAGonStar documentation (\texttt{docs/}) and tutorial sources
(\texttt{docs/tutorial/}, \texttt{examples/}) on the
\texttt{development-native-web} branch of
\url{https://github.com/DagOnStar/dagonstar}. Where the documentation
explicitly flags a limitation, an open technical-debt item, or a constraint of
the current implementation, this book preserves that caveat rather than
smoothing it away, in keeping with the project's own stated principle that
``documentation must remain consistent with the software in this
repository.''
